\begin{abstract}
Sequential LoRA adaptation provides a memory-efficient solution for parameter-efficient continual learning, but a single shared adapter is prone to catastrophic forgetting because its effective update $\Delta W=BA$ drifts away from configurations that support previous tasks. Existing remedies typically address either current-task plasticity through sharpness-aware optimization or old-task stability through Fisher-weighted regularization. A naive combination applies sharpness-aware optimization to the summed objective, but this couples the current-task perturbation with the past-task Fisher penalty and mixes two functionally distinct signals.

We propose \textbf{Flat-Stable LoRA (FS-LoRA)}, a decoupled optimizer for continual LoRA adaptation. FS-LoRA computes an adapter-subspace first-order sharpness correction using only the current-task loss, while adding a Fisher stability gradient as an independent component in the effective adapter-update space $\Delta W=BA$. To prevent the Fisher term from becoming inactive under low-rank parameterization, FS-LoRA further applies trace normalization to the Fisher estimate, making the stability weight scale more consistently across datasets. The method can also be interpreted as a penalty-method approximation to a constrained optimization problem that minimizes current-task loss and sharpness subject to a Fisher drift budget.

Across six fine-grained continual recognition benchmarks, we observe that the flatness correction and the Fisher gradient are nearly orthogonal, with cosine values close to zero throughout training. This gradient-geometric diagnostic supports the decoupled design and explains why flatness and stability should not be mixed in a shared perturbation closure. Empirically, FS-LoRA improves final average accuracy over SeqLoRA with sharpness-aware optimization on several datasets, while revealing mixed behavior in cases where flatness scale or Fisher activity is poorly matched. These results show that separating and scale-calibrating plasticity and stability gradients provides a practical design principle for continual LoRA optimization.
\end{abstract}